\Article{19}
Zerlegt man die magnetische Kraft im Punkt $O$ in drei zueinander senkrechte Komponenten $X,Y,Z$, von denen die dritte zum Erdmittelpunkt gerichtet ist, so liegen $X$ und $Y$ tangential an der durch $O$ verlaufenden, mit der Erde konzentrischen Kugelfläche. Dabei ist $X$ in der durch $O$ und die Erdachse verlaufenden Ebene nach Norden und $Y$ parallel zum Erdäquator nach Westen gerichtet. Dann gilt
\[
X=-\frac{\mathrm{d}V}{r\mathrm{d}u},\quad Y=-\frac{\mathrm{d}V}{r\sin u\,\mathrm{d}\lambda},\quad Z=-\frac{\mathrm{d}V}{\mathrm{d}r}
\]
und folglich
\[
\begin{aligned}
X&=-\frac{R^3}{r^3}\left(\frac{\mathrm{d}P'}{\mathrm{d}u}+\frac Rr\cdot\frac{\mathrm{d}P''}{\mathrm{d}u}+\frac{RR}{rr}\cdot\frac{\mathrm{d}P'''}{\mathrm{d}u}+\text{u. s. w.}\right)\\
Y&=-\frac{R^3}{r^3\sin u}\left(\frac{\mathrm{d}P'}{\mathrm{d}\lambda}+\frac Rr\cdot\frac{\mathrm{d}P''}{\mathrm{d}\lambda}+\frac{RR}{rr}\cdot\frac{\mathrm{d}P'''}{\mathrm{d}\lambda}+\text{u. s. w.}\right)\\
Z&=\frac{R^3}{r^3}\left(2P'+\frac{3RP''}{r}+\frac{4RRP'''}{rr}+\text{u. s. w.}\right)
\end{aligned}
\]

Auf der Erdoberfläche sind $X,Y$ dieselben horizontalen Komponenten, die oben mit diesen Buchstaben bezeichnet wurden, und $Z$ ist die vertikale Komponente, positiv, wenn sie nach unten gerichtet ist. Die Ausdrücke für diese Kräfte auf der Erdoberfläche lauten also
\SourceLeaf{147}
\[
\begin{aligned}
X&=-\left(\frac{\mathrm{d}P'}{\mathrm{d}u}+\frac{\mathrm{d}P''}{\mathrm{d}u}+\frac{\mathrm{d}P'''}{\mathrm{d}u}+\text{u. s. w.}\right)\\
Y&=-\frac1{\sin u}\left(\frac{\mathrm{d}P'}{\mathrm{d}\lambda}+\frac{\mathrm{d}P''}{\mathrm{d}\lambda}+\frac{\mathrm{d}P'''}{\mathrm{d}\lambda}+\text{u. s. w.}\right)\\
Z&=2P'+3P''+4P'''+\text{u. s. w.}
\end{aligned}
\]

\Article{20}
Verbinden wir nun mit diesen Sätzen das bekannte Theorem, dass jede Funktion von $\lambda$ und $u$, die für alle Werte von $\lambda$ von $0$ bis $360^0$ und von $u$ von $0$ bis $180^0$ einen bestimmten endlichen Wert hat, in eine Reihe der Form
\[
P^0+P'+P''+P'''+\text{u. s. w.}
\]
entwickelt werden kann, deren allgemeines Glied $P^{\mathrm{n}}$ die obige partielle Differentialgleichung erfüllt, dass eine solche Entwicklung nur auf eine bestimmte Weise möglich ist und dass diese Reihe stets konvergiert, so erhalten wir folgende bemerkenswerte Sätze.

I. Die Kenntnis des Werts von $V$ an allen Punkten der Erdoberfläche genügt, um daraus den allgemeinen Ausdruck für $V$ im gesamten unendlichen Raum außerhalb der Erdoberfläche abzuleiten und somit auch die Kräfte $X,Y,Z$ auf der Erdoberfläche und im gesamten unendlichen Raum außerhalb davon zu bestimmen. Offenbar braucht man dazu nur $\frac VR$ nach dem erwähnten Theorem in eine Reihe zu entwickeln.

Im Folgenden ist das Zeichen $V$ daher stets in der auf die Erdoberfläche beschränkten Bedeutung zu verstehen, sofern nicht ausdrücklich das Gegenteil gesagt wird, also als diejenige Funktion von $\lambda$ und $u$, die sich aus dem allgemeinen Ausdruck ergibt, wenn man $r=R$ setzt. Damit gilt
\[
V=R(P'+P''+P'''+\text{u. s. w.})
\]

II. Die Kenntnis des Werts von $X$ an allen Punkten der Erdoberfläche genügt, um alles in I. Genannte zu erhalten. Tatsächlich gilt nach Art. 15 für das Integral $\int_0^u X\mathrm{d}u=\frac{V^0-V}{R}$, wobei $V^0$ den Wert von $V$ am Nordpol bezeichnet. Die Entwicklung von $\int_0^u X\mathrm{d}u$ in eine Reihe der erwähnten Form muss mit
\SourceLeaf{148}
\[
V^0-P'-P''-P'''-\text{u. s. w.}
\]
identisch sein.

III. Ebenso ist mit Bezug auf Art. 16 klar, dass die Kenntnis des Werts von $Y$ auf der ganzen Erde zusammen mit der Kenntnis von $X$ an allen Punkten einer von einem Erdpol zum anderen verlaufenden Linie genügt, um die \textit{vollständige} Theorie des Erdmagnetismus zu begründen.

IV. Schließlich ist klar, dass sich die vollständige Theorie auch allein aus der Kenntnis des Werts von $Z$ auf der gesamten Erdoberfläche ableiten lässt. Entwickelt man nämlich $Z$ in eine Reihe
\[
Z=Q^0+Q'+Q''+Q'''+\text{u. s. w.}
\]
deren allgemeines Glied die mehrfach erwähnte partielle Differentialgleichung erfüllt, so muss notwendigerweise $Q^0=0$ gelten und $P'=\frac12Q'$, $P''=\frac13Q''$, $P'''=\frac14Q'''$ u. s. w sein.

\Article{21}
Weil die einzelnen Kräfte $X,Y,Z$ auf so einfache Weise von einer einzigen Funktion $V$ abhängen und auch untereinander in einem einfachen Zusammenhang stehen, eignen sie sich viel besser als Grundlage der Theorie als die übliche Darstellung der magnetischen Kraft durch die drei Elemente Gesamtintensität, Inklination und Deklination. Oder vielmehr, die letztere Darstellungsart. so natürlich sie an sich erscheint, wenn es nur darauf ankommt, die Tatsachen zu erfassen, kann unmittelbar überhaupt nicht zur Begründung der Theorie, wenigstens nicht zu ihrer ersten Begründung, dienen, bevor sie in die andere Form übertragen worden ist. In dieser Hinsicht wäre daher eine allgemeine grafische Darstellung der horizontalen Intensität sehr wünschenswert. Zum einen kommt diese dem, was für die Theorie brauchbar ist, näher als die Gesamtintensität. Zum anderen ist sie in der weit überwiegenden Zahl der Fälle das, was ursprünglich tatsächlich beobachtet wurde, während die Gesamtintensität nur durch Rechnung mithilfe der Inklination daraus abgeleitet wurde. Die Elemente des horizontalen Magnetismus gesondert zu erhalten, ist umso mehr zu empfehlen, als sie sich mit den gegenwärtigen Hilfsmitteln mit deutlich größerer Genauigkeit bestimmen lassen. Zumindest sollte man niemals die aus der beobachteten horizontalen \SourceLeaf{149}Intensität berechnete Gesamtintensität unter Weglassung des beobachteten Werts bekannt geben, ohne die bei der Rechnung verwendete Inklination mit anzugeben, damit diejenigen, die die Daten für die Theorie nutzen wollen, die ursprünglichen Zahlen unverfälscht wiederherstellen können.

So interessant es auch wäre, die ganze Theorie des Erdmagnetismus allein auf Beobachtungen der horizontalen Nadel zu gründen und damit die vertikale Komponente oder die Inklination vorwegzunehmen, so ist es dafür gegenwärtig doch noch viel zu früh. Die Unzulänglichkeit der jetzt verfügbaren Daten erlaubt es nicht, auf die Mitverwendung der vertikalen Komponente zu verzichten. Im Grunde erfährt die Theorie auch schon dadurch eine Bestätigung, dass sich die Vereinbarkeit sämtlicher Elemente unter einem einzigen Prinzip nachweisen lässt.


